Scientific workflows orchestrate complex data analyses on large, compute-intensive datasets that often exceed the capacity of a single computer or even a single cluster.
Offloading jobs across multiple clusters can alleviate resource constraints and shorten execution time while controlling cost. 

This paper proposes adaptive offloading of workflow jobs in multi-cluster environments to meet user-defined makespan deadlines while minimizing execution costs. 
We propose a heuristic based on a two-step prediction approach: (i) predicting individual job runtimes using a regression model trained on historical executions and (ii) simulating workflow execution to estimate makespan and cost for alternative offloading configurations.

The method is implemented as an extension of the Snakemake workflow management system and is evaluated with three real-world bioinformatic workflows. 
By determining the locally optimal subset of jobs to offload, our approach reduces the overall workflow makespan by up to 36.51\% compared to a single-cluster baseline and reduces costs under deadline constraints by approximately 52\% compared to naive offloading, while maintaining absolute makespan prediction errors between 0.21\% and 16.88\%.
